\documentclass[12pt, a4paper]{article}

\usepackage{amsmath, amssymb, amsthm}
\usepackage{geometry}
\usepackage{booktabs}
\usepackage{array}
\usepackage{hyperref}
\usepackage{xcolor}
\usepackage{tcolorbox}
\usepackage{enumitem}
\usepackage{parskip}
\usepackage{tikz}
\usetikzlibrary{arrows.meta, positioning, calc}

\geometry{margin=2.5cm}

\newif\ifstudentmode
\studentmodefalse
\ifx\studentflag\undefined\else\studentmodetrue\fi

\newcommand{\sol}[2][6cm]{%
  \ifstudentmode\underline{\hspace{#1}}\else{#2}\fi}
\newcommand{\soleq}[2][3cm]{%
  \ifstudentmode
    \par\vspace{4pt}\noindent\rule{\linewidth}{0.5pt}\\[#1]%
    \noindent\rule{\linewidth}{0.5pt}\par\vspace{4pt}%
  \else\[#2\]\fi}

\definecolor{mainblue}{RGB}{30, 100, 180}
\definecolor{boxbg}{RGB}{235, 243, 255}
\definecolor{notebg}{RGB}{255, 248, 230}
\definecolor{warnbg}{RGB}{255, 235, 235}
\definecolor{greenbg}{RGB}{230, 255, 235}

\tcbuselibrary{skins, breakable}

\newtcolorbox{resultbox}{colback=boxbg,colframe=mainblue,boxrule=1.2pt,arc=4pt,
  title={Result},fonttitle=\bfseries\color{white},coltitle=white,
  attach boxed title to top left={yshift=-2mm,xshift=4mm},
  boxed title style={colback=mainblue,arc=3pt}}

\newtcolorbox{notebox}{colback=notebg,colframe=orange!70!black,
  boxrule=1pt,arc=4pt,left=4pt,right=4pt}

\newtcolorbox{warnbox}{colback=warnbg,colframe=red!60!black,
  boxrule=1pt,arc=4pt,left=4pt,right=4pt}

\newtcolorbox{stepbox}[1]{colback=greenbg,colframe=green!50!black,
  boxrule=1pt,arc=4pt,title={#1},fonttitle=\bfseries,coltitle=green!30!black}

\theoremstyle{definition}
\newtheorem{example}{Example}[section]

\newcommand{\muB}{\mu_{\mathcal{B}}}
\newcommand{\sigB}{\sigma^2_{\mathcal{B}}}

\begin{document}

\begin{center}
  {\LARGE \textbf{Section 4.4: Batch Normalisation}} \\[0.3em]
  {\large\itshape Keeping Activations Well-Behaved During Training} \\[0.5em]
  {\large AI Class --- Lecture Notes}%
  \ifstudentmode\quad{\normalsize\color{gray}[Student Copy]}\fi
\end{center}

\vspace{1em}\hrule\vspace{1.5em}

% ══════════════════════════════════════════════════════════════════════════════
\section{The Problem: Shifting Activations}

As weights update during training, the distribution of pre-activations
entering each layer shifts.  The next layer must constantly adapt to a
moving input --- slowing training and making it unstable.

For example, if the pre-activations $z$ drift to large positive values,
the sigmoid $\sigma(z) \to 1$ and its gradient $\sigma'(z) \to 0$.
Gradients vanish and training stalls.

\textbf{Batch normalisation (BN)} fixes this by normalising the
pre-activations to zero mean and unit variance \emph{within each
mini-batch}, right before the activation function.

% ══════════════════════════════════════════════════════════════════════════════
\section{The Algorithm}

Given a mini-batch $\mathcal{B} = \{z_1, z_2, \ldots, z_m\}$ of
pre-activations for one neuron:

\begin{resultbox}
\textbf{Step 1.} Compute the batch mean and variance:
\[
  \muB = \frac{1}{m}\sum_{i=1}^{m} z_i,
  \qquad
  \sigB = \frac{1}{m}\sum_{i=1}^{m}(z_i - \muB)^2
\]

\textbf{Step 2.} Normalise each value:
\[
  \hat{z}_i = \frac{z_i - \muB}{\sqrt{\sigB + \varepsilon}}
\]

\textbf{Step 3.} Scale and shift with learned parameters $\gamma$ and $\beta$:
\[
  \tilde{z}_i = \gamma\,\hat{z}_i + \beta
\]

The output $\tilde{z}_i$ is passed to the activation function.
\end{resultbox}

$\varepsilon \approx 10^{-5}$ prevents division by zero.
$\gamma$ and $\beta$ are \textbf{learnable}: the network can undo the
normalisation if that turns out to be best for the task.

% ══════════════════════════════════════════════════════════════════════════════
\section{Hand-Calculated Example}

Mini-batch of $m = 4$ pre-activations (before the hidden-layer sigmoid):
\[
  \mathcal{B} = \{0.5,\; 1.2,\; -0.3,\; 0.8\}
\]

% ──────────────────────────────────────────────────────────────────────────────
\begin{stepbox}{Step 1 \quad Batch mean and variance}

\[
  \muB = \frac{0.5 + 1.2 + (-0.3) + 0.8}{4}
       = \frac{2.2}{4}
       = \sol[3cm]{0.55}
\]

Deviations $(z_i - \muB)$:

\begin{center}
\renewcommand{\arraystretch}{1.3}
\begin{tabular}{cccc}
\toprule
$z_i$ & $z_i - \muB$ & $(z_i - \muB)^2$ \\
\midrule
$0.5$  & $-0.05$ & $0.0025$ \\
$1.2$  & $\phantom{-}0.65$ & $0.4225$ \\
$-0.3$ & $-0.85$ & $0.7225$ \\
$0.8$  & $\phantom{-}0.25$ & $0.0625$ \\
\midrule
       &         & $\Sigma = 1.2100$ \\
\bottomrule
\end{tabular}
\end{center}

\[
  \sigB = \frac{1.2100}{4} = \sol[3cm]{0.3025}
  \qquad
  \sqrt{\sigB} = \sol[3cm]{0.55}
\]
\end{stepbox}

\vspace{0.8em}

% ──────────────────────────────────────────────────────────────────────────────
\begin{stepbox}{Step 2 \quad Normalise}

Divide each deviation by $\sqrt{\sigB} = 0.55$ (ignoring $\varepsilon$):

\begin{center}
\renewcommand{\arraystretch}{1.4}
\begin{tabular}{ccc}
\toprule
$z_i$ & $z_i - \muB$ & $\hat{z}_i = (z_i-\muB)\,/\,0.55$ \\
\midrule
$0.5$  & $-0.05$ & \sol[3cm]{$-0.091$} \\
$1.2$  & $\phantom{-}0.65$ & \sol[3cm]{$\phantom{-}1.182$} \\
$-0.3$ & $-0.85$ & \sol[3cm]{$-1.545$} \\
$0.8$  & $\phantom{-}0.25$ & \sol[3cm]{$\phantom{-}0.455$} \\
\bottomrule
\end{tabular}
\end{center}

Verify: mean$(\hat{z}) \approx 0$ and variance$(\hat{z}) \approx 1$. \checkmark
\end{stepbox}

\vspace{0.8em}

% ──────────────────────────────────────────────────────────────────────────────
\begin{stepbox}{Step 3 \quad Scale and shift}

Using the initial values $\gamma = 1$, $\beta = 0$
(the network will learn these):
\[
  \tilde{z}_i = 1 \cdot \hat{z}_i + 0 = \hat{z}_i
\]

The output $\tilde{z}_i$ is passed to the activation function.  With
$\gamma = 1, \beta = 0$, BN simply standardises the batch.
After training, $\gamma$ and $\beta$ may take any value.
\end{stepbox}

% ══════════════════════════════════════════════════════════════════════════════
\section{Where to Place Batch Normalisation}

BN is inserted \textbf{between the linear transformation and the activation
function}:

\[
  \mathbf{z} = \mathbf{W}\mathbf{x} + \mathbf{b}
  \;\longrightarrow\;
  \tilde{\mathbf{z}} = \text{BN}(\mathbf{z})
  \;\longrightarrow\;
  \mathbf{h} = \sigma(\tilde{\mathbf{z}})
\]

\begin{notebox}
When BN is used, the bias $\mathbf{b}$ in the linear layer is redundant:
BN subtracts the batch mean, which removes any constant shift.
In PyTorch, set \texttt{bias=False} in a linear layer that is followed by
\texttt{nn.BatchNorm1d}.
\end{notebox}

% ══════════════════════════════════════════════════════════════════════════════
\section{Training vs Inference}

\begin{center}
\renewcommand{\arraystretch}{1.3}
\begin{tabular}{lll}
\toprule
Phase & Mean and variance used & Source \\
\midrule
Training   & Mini-batch $\muB$, $\sigB$ & Computed from current batch \\
Inference  & Running mean $\bar\mu$, running var $\bar\sigma^2$ &
             Accumulated from all training batches \\
\bottomrule
\end{tabular}
\end{center}

During training, PyTorch tracks a running average of $\muB$ and $\sigB$
over all batches.  At inference, these stored statistics are used so that
predictions on single examples are stable.

\begin{warnbox}
Like dropout, BN behaves differently in training and inference mode.
Always call \texttt{model.train()} before training and
\texttt{model.eval()} before evaluation.
Forgetting this is one of the most common bugs in PyTorch projects.
\end{warnbox}

\begin{notebox}
\textbf{PyTorch.}
\texttt{nn.BatchNorm1d(num\_features)} for fully-connected layers;
\texttt{nn.BatchNorm2d(num\_channels)} for convolutional layers.
See \texttt{code\_examples.py} for a complete example.
\end{notebox}

\end{document}
